\cleardoublepage
\section{Inequalities}

In this section, we'll learn how to solve basic inequalities and how to graph your answers on a number line. Inequalities use signs like:
\begin{itemize}
  \item $>$ greater than
  \item $<$ less than
  \item $\geq$ greater than or equal to
  \item $\leq$ less than or equal to
\end{itemize}

Solve inequalities just like equations, \textbf{but remember}: If you multiply or divide both sides by a \textbf{negative number}, you must \textbf{flip the inequality sign}.

Inequalities are used in many jobs—like making decisions in business, writing computer programs, or planning safe speeds in transportation.

\subsection*{Example: Solving Inequalities}

Let's look at some examples:

\textbf{Basic Inequality:}
\[
x - 3 > 2 \quad \text{Add 3 to both sides} \quad x > 5
\]

\textbf{Inequality with Negative Coefficient:}
\[
-2x < 6 \quad \text{Divide both sides by -2, flip the sign} \quad x > -3
\]

Suppose a sign says, "You must be \textbf{under 18} to enter." That's an inequality! It means your age must be \textbf{less than 18}, or $x < 18$. Inequalities describe rules, limits, and ranges all around us.

Always treat inequalities like equations—unless you multiply or divide by a negative! Graph solutions using number lines. Use:
\begin{itemize}
  \item \textbf{Open circles} for $<$ or $>$
  \item \textbf{Closed circles} for $\leq$ or $\geq$
\end{itemize}

Check your solution by testing a number that fits the inequality.

\cleardoublepage
\subsection{Short Question (English)}
Solve: $x + 4 < 10$

\fullpageimage{images/q2_5_1-eng.jpg}

\newpage
\subsection{Short Question (Korean)}
Solve: $x + 4 < 10$

\fullpageimage{images/q2_5_1-kor.jpg}

\newpage
\subsection{Short Question (Answer)}
Solve: $x + 4 < 10$

\textbf{Answer:} $x < 6$

\cleardoublepage
\subsection{Long Question (English)}
Solve and graph: $-3x \geq 12$

\fullpageimage{images/q2_5_2-eng.jpg}

\newpage
\subsection{Long Question (Korean)}
Solve and graph: $-3x \geq 12$

\fullpageimage{images/q2_5_2-kor.jpg}

\newpage
\subsection{Long Question (Answer)}
Solve and graph: $-3x \geq 12$

\textbf{Answer:}\\
Divide by $-3$ and flip the sign: $x \leq -4$\\
Graph: closed circle at –4, shade to the left.

\cleardoublepage
\subsection{Challenge Question (English)}
A movie ticket costs \$12. Write and solve an inequality for the number of tickets you can buy with \$80.

\fullpageimage{images/q2_5_3-eng.jpg}

\newpage
\subsection{Challenge Question (Korean)}
A movie ticket costs \$12. Write and solve an inequality for the number of tickets you can buy with \$80.

\fullpageimage{images/q2_5_3-kor.jpg}

\newpage
\subsection{Challenge Question (Answer)}
A movie ticket costs \$12. Write and solve an inequality for the number of tickets you can buy with \$80.

\textbf{Answer:}\\
$12t \leq 80$\\
$t \leq 6.\overline{6}$ → You can buy at most \textbf{6 tickets}.

